\begin{theorem}[Empty quadratic common-line image]
\label{thm:fixed-face-quadratic-empty}
Over \(\mathbb F_5\) the face \([m_2:m_3]=[0:1]\) has empty
quadratic common-line incidence. On its eligible quotient \(U\),
\[
             Z=Z_{\rm lev}=R_{\rm lev}=0 .
\]
In particular the residual has degree zero and its support contains
zero Frobenius orbits.
\end{theorem}

\begin{proof}
Write \(k=\mathbb F_5\). Use the three quadrics of the net in turn
to test the same polar cofactor \(c(x)\). For \(i=2,3,4\) put
\[
  B_D^2 G_i=\operatorname{Res}_x(H,Q_2,Q_3,Q_4,Q_i(c(x))),
  \qquad d_i=336+8(i-2).
\]
Proposition~\ref{prop:common-line-residual} constructs \(G_2\).
Its auxiliary-square proof applies to all three forms. At a generic
section point \(x\), write the polar kernel as \(\langle x,y\rangle\).
Since \(Q_i(x)=b_i(x,y)=0\), evaluating \(Q_i\) on the cofactor
produces the factor \(D(x)^2\) for each of the eight points. The same
universal polynomial division and cancellation on the framed incidence
give \(G_i=0\) on every common line, including specializations where
\(B_D=0\). Replacing \(Q_2(c)\) by \(Q_i(c)\) adds the coefficient
twist \(\mathcal O_C(i-2)\) to the sextic, whose resultant degree is
eight. Thus \(G_i\) is an invariant section of \(\mathcal O_C(d_i)\).
Every possible common line is covered by \(Y\ne0\) by
Theorem~\ref{thm:common-line-decomposition}. There we use
\(g_i=G_i/Y^{d_i}\).

\emph{Exact section values.}
Write \(x_0,x_1,x_2\) for the edges incident to vertex 2 and \(y_j\)
for the opposite edge paired with \(x_j\). Then \(\sum x_j=0\) and
\(Y=\sum y_j\). The calculation uses the affine frame \(Y=1\)
and the invariant parameter \(t=x_0x_1x_2\) in this frame.
At each supplied pure point, expand
\(\operatorname{per}(J_4+sB+uX)\) to obtain the three polar rows
and the three Gram matrices. In the nine-dimensional tangent space,
let \(M\) have rows consisting of the three polars, a coordinate
functional \(\ell\) with \(\ell(B)\ne0\), three more coordinate rows,
and \(D,H\), in that order. Choose the three rows so \(M\) is
invertible. The last five columns of \(M^{-1}\) give an adapted
frame of \(W_B\); the volume in this frame is
\(\pm\ell(B)/\det M\) by expansion along the polar and gauge rows.

In the adapted frame let \(\beta\) be the ternary resultant of the
quadrics at \(D=H=0\). At every sample it is nonzero, so \(H=0,D=1\)
contains the whole length-eight section with affine algebra \(A_B\).
The homogeneous resultant product formula, with multiplicities
\cite{jouanolouResultant},
gives the following scalar in the fixed ambient trivialization:
\[
   g_i(B)=\left(\frac{\det M}{\ell(B)}\right)^{48}
             \beta^4\,\det\!\left(m_{Q_i(c)}:A_B\to A_B\right).
\]
Indeed the sextic gives \(\beta^6\) before canceling \(\beta^2\).
The determinant twist of \(G_i\) is \((\det W^*)^{48}\) for every
\(i\). The evaluator computes \(\beta\) by the ternary Macaulay
determinant quotient and reduces multiples of the three quadrics
through degree four to obtain the eight-dimensional multiplication
matrices.

\emph{Certification on the entire curve.}
Use the field \(k[w]/(w^7+3w+3)\) for the sample evaluations.
The checker substitutes all 450 labeled points in the pure equations.
Each has three distinct incident edges. Permuting the three edge pairs
and applying Frobenius gives \(450\cdot6\cdot7=18\,900\) distinct
geometric points of \(C\), as checked on the normalized coordinates.
The resultant evaluation also agrees in two frames at each of the
six labelings of every sample.
Fit each \(g_i\) by an invariant polynomial \(P_i\) using the columns
\(t^j b_h\), where the \(b_h\) are orbit sums of monomials in
\(x_0,x_1,y_0,y_1\), \(b_0=1\), and \(3j+\deg b_h\le d_i\).
The twelve orbit sums have degrees
\[
              0,2,2,3,3,4,4,5,5,6,6,8 .
\]
They are obtained by row reduction of the orbit sums modulo the pure
equations and \(Y=x_0x_1x_2=0\). The checker verifies the weight of
every recovered term and all 450 equalities in the extension field.
Homogenize \(P_i\) with \(Y\) to degree \(d_i\). Its difference
from \(G_i\) vanishes on all the permuted and conjugate points by
invariance and descent to \(k\). A nonzero section of
\(\mathcal O_C(d_i)\) has zero divisor of degree \(24d_i\)
\cite{stacksCurveDegrees}, since
\(C\) is geometrically integral of degree 24. Even the largest
bound is \(24\cdot352=8\,448\). Thus each recovered polynomial
equals the true section on \(Y=1\).

\emph{Exclusion by two coprime determinants.}
In the coordinate ring of this chart, reduce the 78 products
\(b_h b_j\), \(h\le j\), modulo the pure equations. The resulting
identities express each product in the \(b_h\) with coefficients in
\(k[t]\). They give a \(12\times12\) polynomial matrix \(M_i(t)\)
for multiplication by \(P_i\), and \(N_i(t)=\det M_i(t)\).
In particular, at any geometric base point \(s\), evaluation gives
\[
 (b_h(s))_h\,M_i(t(s))=g_i(s)\,(b_h(s))_h .
\]
This row is nonzero because \(b_0=1\). A common line above \(s\)
would force all three \(N_i(t(s))\) to vanish. Exact polynomial
arithmetic over \(k\) gives
\[
  (\deg N_2,\deg N_3,\deg N_4)=(1344,1372,1406),
  \qquad \gcd(N_2,N_3)=\gcd(N_2,N_4)=\gcd(N_3,N_4)=1 .
\]
Already the first pair excludes every value of \(t\) over
\(\overline k\). Consequently \(\mathcal F\) and its labeled
pullback have no geometric points and are empty schemes. Their
images and the residual divisor are zero.
\end{proof}

\paragraph{Finiteness of kernel directions.}
The entire projective locus of pure kernel directions admitting a
quadratic common line is finite over \(\overline{\mathbb F}_5\) and
over \(\mathbb C\). Indeed the theorem excludes \(m_2=0,m_3\ne0\)
at five, and Proposition~\ref{prop:kernel-curve} excludes
\(m_2=m_3=0\). The projective flag incidence used there, now with
the single equation \(m_2=0\), has empty special fiber over
\(\mathbb Z_{(5)}\). Properness gives an empty generic fiber as well.
In both characteristics the closed direction image therefore misses
a hyperplane in \(\mathbb P^5\). A projective scheme contained in
this affine complement is finite over its field. By symmetry any
surviving direction has all four vertex sums nonzero. Existence
of such a direction off the hyperplanes remains open.
